%=============================================================================
% PSI-BUBBLE OBJECT CAPTURE AND FREEFALL DYNAMICS
% Derived from DFD Formalism (Section 2 of the Unified Theory)
% Agent 2: Missile Capture, Tidal Forces, EM Propagation, Atmospheric Transit
%=============================================================================

\documentclass[12pt,aps,prd,superscriptaddress,nofootinbib]{revtex4-2}
\usepackage{amsmath,amsfonts,amssymb,bm,graphicx,siunitx,physics}

\newcommand{\psifield}{\psi}
\newcommand{\avec}{\mathbf{a}}
\newcommand{\neff}{n_{\rm eff}}
\newcommand{\psib}{\psi_{\rm b}}
\newcommand{\Rb}{R_{\rm b}}
\newcommand{\Rc}{R_{\rm c}}
\newcommand{\Lshell}{\ell_{\rm shell}}

\begin{document}

\title{Psi-Bubble Object Capture via Freefall:\\
Missile Entrainment, Tidal Forces, EM Refraction,\\
and Atmospheric Transit in DFD}

\author{DFD Research Programme}
\date{\today}

\begin{abstract}
We derive, from the DFD formalism, the complete physics of object
capture by a technologically generated $\psi$-gradient bubble.
A device producing an asymmetric $\neff = \exp[\psi + \kappa(\eta - \eta_c)]$
distribution creates a local gravitational well.
Any object entering the gradient region---including projectiles,
debris, or atmospheric molecules---falls freely in the same
$\psi$-gradient as the ship, experiencing \emph{zero relative
acceleration} (the Weak Equivalence Principle, DFD Eq.~2.14).
We compute: (1)~the capture radius where the $\psi$-gradient
becomes dynamically significant; (2)~tidal forces at the bubble
boundary; (3)~EM signal refraction and lensing through the
$n \neq 1$ region; (4)~external observables (spectral shifts,
radar returns, luminosity changes); and (5)~atmospheric transit
dynamics.
\end{abstract}

\maketitle
\tableofcontents

%=============================================================================
\section{Foundational Framework}\label{sec:foundations}
%=============================================================================

\subsection{The DFD Equation of Motion}

From the DFD formalism (Section~2, Eq.~2.14), the acceleration of
\emph{any} test mass in a $\psi$-field is:
\begin{equation}
\frac{d^2\mathbf{x}}{dt^2} = \frac{c^2}{2}\nabla\psi \equiv \avec.
\label{eq:eom}
\end{equation}
This is \textbf{mass-independent}: a missile of mass $m_{\rm missile}$
and a ship of mass $m_{\rm ship}$ in the same $\psi$-gradient
experience identical acceleration. This is the DFD realization of the
Weak Equivalence Principle---gravity is universal freefall in the
refractive medium.

\subsection{The Engineered $\psi$-Bubble}

Consider a device that generates a modified effective refractive index
(DFD Section~11A, Eq.~11.4):
\begin{equation}
\neff(\mathbf{x}) = \exp\!\Big[\psi_{\rm bg}(\mathbf{x})
  + \psi_{\rm eng}(\mathbf{x})\Big],
\label{eq:neff}
\end{equation}
where $\psi_{\rm bg}$ is the ambient gravitational $\psi$-field
(e.g., Earth's gravity) and $\psi_{\rm eng}$ is the engineered
contribution. For propulsion, the device creates an
\textbf{asymmetric} $\psi_{\rm eng}$: larger in the direction of
desired travel, producing a net gradient.

We model the engineered field as a shell-like distribution with
characteristic parameters:
\begin{itemize}
\item $\psib$: peak engineered $\psi$-value at the shell boundary
\item $\Rb$: radius of the bubble (distance from device to shell peak)
\item $\Lshell$: thickness of the gradient transition layer
\item $\hat{\mathbf{d}}$: direction of asymmetry (direction of travel)
\end{itemize}

For the symmetric component, we adopt a smooth profile:
\begin{equation}
\psi_{\rm eng}(r) = \psib\,\operatorname{sech}^2\!\left(
  \frac{r - \Rb}{\Lshell}\right),
\label{eq:psi-profile}
\end{equation}
which peaks at $r = \Rb$ and decays on scale $\Lshell$ both inward
and outward. The gradient is:
\begin{equation}
\frac{\partial\psi_{\rm eng}}{\partial r} =
-\frac{2\psib}{\Lshell}\,\operatorname{sech}^2\!\left(
  \frac{r - \Rb}{\Lshell}\right)
\tanh\!\left(\frac{r - \Rb}{\Lshell}\right).
\label{eq:psi-gradient}
\end{equation}
The maximum gradient magnitude occurs at $r = \Rb \pm \Lshell/2$
and equals:
\begin{equation}
|\nabla\psi|_{\rm max} \approx \frac{\psib}{\Lshell}.
\label{eq:grad-max}
\end{equation}

The corresponding maximum acceleration (from Eq.~\ref{eq:eom}) is:
\begin{equation}
\boxed{a_{\rm max} = \frac{c^2}{2}\frac{\psib}{\Lshell}.}
\label{eq:a-max}
\end{equation}

\paragraph{Numerical scale.}
For $a_{\rm max} = 100g \approx 10^3\;\si{m/s^2}$ and
$\Lshell = 10\;\si{m}$:
\begin{equation}
\psib = \frac{2\,a_{\rm max}\,\Lshell}{c^2}
= \frac{2 \times 10^3 \times 10}{(3\times 10^8)^2}
\approx 2.2 \times 10^{-13}.
\end{equation}
This is a minuscule $\psi$-perturbation---yet it produces
macroscopic gravitational acceleration.


%=============================================================================
\section{Capture Radius}\label{sec:capture}
%=============================================================================

\subsection{Definition}

The \textbf{capture radius} $\Rc$ is the distance from the bubble
centre at which the engineered $\psi$-gradient produces an
acceleration exceeding a threshold $a_{\rm thresh}$:
\begin{equation}
a(r) = \frac{c^2}{2}\left|\frac{\partial\psi_{\rm eng}}
  {\partial r}\right| \geq a_{\rm thresh}.
\label{eq:capture-def}
\end{equation}

\subsection{Computation}

Using the $\operatorname{sech}^2\tanh$ profile
(Eq.~\ref{eq:psi-gradient}), the acceleration at distance $r$ from
bubble centre is:
\begin{equation}
a(r) = \frac{c^2\psib}{\Lshell}\,\operatorname{sech}^2\!\left(
  \frac{r - \Rb}{\Lshell}\right)
\left|\tanh\!\left(\frac{r - \Rb}{\Lshell}\right)\right|.
\end{equation}

In the far field ($r \gg \Rb + \Lshell$), the $\operatorname{sech}^2$
factor dominates and we have:
\begin{equation}
a(r) \approx 4\,a_{\rm max}\,
  \exp\!\left(-\frac{2(r - \Rb)}{\Lshell}\right),
  \qquad r \gg \Rb.
\end{equation}

Setting $a(r) = a_{\rm thresh}$ and solving:
\begin{equation}
\boxed{\Rc = \Rb + \frac{\Lshell}{2}\ln\!\left(
  \frac{4\,a_{\rm max}}{a_{\rm thresh}}\right).}
\label{eq:capture-radius}
\end{equation}

\paragraph{Numerical example.}
For $a_{\rm max} = 100g$, $\Rb = 50\;\si{m}$,
$\Lshell = 10\;\si{m}$, and $a_{\rm thresh} = 1g$:
\begin{equation}
\Rc = 50 + 5\,\ln(400) \approx 50 + 30 = 80\;\si{m}.
\end{equation}
The capture zone extends roughly $30\;\si{m}$ beyond the nominal
bubble shell. For $a_{\rm thresh} = 0.01g$ (detectable deflection):
\begin{equation}
\Rc \approx 50 + 5\,\ln(40000) \approx 103\;\si{m}.
\end{equation}

\subsection{Scaling Relations}

The capture radius depends logarithmically on the acceleration ratio,
but linearly on shell thickness:
\begin{equation}
\Rc - \Rb \propto \Lshell\,\ln\!\left(
  \frac{a_{\rm max}}{a_{\rm thresh}}\right).
\end{equation}

A thicker gradient shell (larger $\Lshell$) produces a larger
capture zone. A steeper bubble (smaller $\Lshell$ for fixed $\psib$)
produces a more intense but more compact capture region.


%=============================================================================
\section{Missile Capture: The Freefall Mechanism}\label{sec:missile}
%=============================================================================

\subsection{Setup}

A missile of mass $m$ approaches the $\psi$-bubble at velocity
$\mathbf{v}_{\rm rel}$ relative to the ship. At time $t_0$, it
crosses the capture radius $\Rc$.

\subsection{Equations of Motion}

In the ship's reference frame, the missile satisfies:
\begin{equation}
\frac{d^2\mathbf{x}_{\rm missile}}{dt^2}
= \frac{c^2}{2}\nabla\psi(\mathbf{x}_{\rm missile})
+ \mathbf{a}_{\rm missile,\,propulsive},
\end{equation}
while the ship (and everything attached to it within the bubble)
satisfies:
\begin{equation}
\frac{d^2\mathbf{x}_{\rm ship}}{dt^2}
= \frac{c^2}{2}\nabla\psi(\mathbf{x}_{\rm ship}).
\end{equation}

The \textbf{relative acceleration} is:
\begin{equation}
\mathbf{a}_{\rm rel} = \frac{c^2}{2}\Big[
  \nabla\psi(\mathbf{x}_{\rm missile})
  - \nabla\psi(\mathbf{x}_{\rm ship})\Big]
  + \mathbf{a}_{\rm missile,\,propulsive}.
\end{equation}

\subsection{The Key Result: Freefall Equivalence}

If the missile engine cuts out (or its thrust is negligible compared
to the $\psi$-gradient acceleration), and if the missile is close
enough to the ship that $\nabla\psi$ is approximately uniform
(i.e., within the tidal tolerance zone), then:
\begin{equation}
\boxed{\mathbf{a}_{\rm rel} \approx \frac{c^2}{2}\,
  \Delta\mathbf{x}\cdot\nabla\nabla\psi = \text{tidal terms only}.}
\label{eq:freefall}
\end{equation}

\textbf{The missile falls freely in the same $\psi$-gradient as
the ship.} It is captured and co-moves with the bubble, exactly as
a rock dropped inside an orbiting spacecraft floats alongside
the astronauts.

\subsection{Capture Dynamics}

\begin{enumerate}
\item \textbf{Approach} ($r > \Rc$): Missile travels on its initial
trajectory, unaffected by the bubble.

\item \textbf{Gradient encounter} ($r \sim \Rc$): Missile enters the
gradient zone. It begins to accelerate at
$a \sim (c^2/2)\,|\nabla\psi_{\rm eng}|$.

\item \textbf{Deceleration of relative motion}: The $\psi$-gradient
acts on the missile in the same direction as the ship's motion.
The missile's velocity relative to the ship decreases as both
are accelerated by the same field.

\item \textbf{Capture} ($r < \Rb$): Inside the bubble, the
$\psi$-field is approximately uniform. The missile floats in
freefall with essentially zero relative acceleration.

\item \textbf{Co-motion}: From this point, the missile moves with
the bubble. From an external observer, it appears to have been
``grabbed'' and accelerated to the bubble's velocity.
\end{enumerate}

\subsection{Capture Timescale}

A missile approaching with relative velocity $v_{\rm rel}$ through
a gradient zone of thickness $\Lshell$ with peak acceleration
$a_{\rm max}$ experiences a velocity change:
\begin{equation}
\Delta v \sim a_{\rm max}\,\frac{\Lshell}{v_{\rm rel}}
\sim \frac{c^2\psib}{2v_{\rm rel}}.
\end{equation}
Capture occurs when $\Delta v \gtrsim v_{\rm rel}$, i.e.:
\begin{equation}
v_{\rm rel} \lesssim v_{\rm capture}
= \sqrt{\frac{c^2\psib}{2}} = c\sqrt{\frac{\psib}{2}}.
\end{equation}
For $\psib \sim 10^{-13}$:
$v_{\rm capture} \sim 3\times 10^8 \times 7\times 10^{-7}
\approx 200\;\si{m/s}$.

For a missile at Mach~3 ($\sim 1000\;\si{m/s}$), one needs
$\psib \sim 2\times 10^{-11}$, corresponding to
$a_{\rm max} \sim 3\times 10^4\;g$. For a slower projectile
(Mach~0.5, $170\;\si{m/s}$), the baseline parameters suffice.


%=============================================================================
\section{Tidal Forces at the Bubble Boundary}\label{sec:tidal}
%=============================================================================

\subsection{Tidal Tensor}

The tidal acceleration across an object of size $L$ within the
$\psi$-gradient is:
\begin{equation}
\Delta a_{\rm tidal} = \frac{c^2}{2}\,|\nabla\nabla\psi|\,L
= \frac{c^2}{2}\left|\frac{\partial^2\psi_{\rm eng}}
  {\partial r^2}\right| L.
\label{eq:tidal}
\end{equation}

\subsection{Computation from the Profile}

Differentiating Eq.~(\ref{eq:psi-gradient}):
\begin{equation}
\frac{\partial^2\psi_{\rm eng}}{\partial r^2}
= \frac{2\psib}{\Lshell^2}\,\operatorname{sech}^2\!\left(
  \frac{r-\Rb}{\Lshell}\right)\left[
  3\tanh^2\!\left(\frac{r-\Rb}{\Lshell}\right) - 1\right].
\end{equation}

The maximum of $|\partial^2\psi/\partial r^2|$ occurs at the
inflection points of the gradient profile, giving:
\begin{equation}
\left|\frac{\partial^2\psi_{\rm eng}}{\partial r^2}\right|_{\rm max}
\sim \frac{2\psib}{\Lshell^2}.
\end{equation}

Therefore the maximum tidal acceleration across an object of size $L$
at the bubble boundary is:
\begin{equation}
\boxed{\Delta a_{\rm tidal,\,max}
  = \frac{c^2\psib}{\Lshell^2}\,L
  = \frac{2\,a_{\rm max}}{\Lshell}\,L.}
\label{eq:tidal-max}
\end{equation}

\subsection{Survivability Assessment}

For a missile of length $L = 5\;\si{m}$, $a_{\rm max} = 100g$,
$\Lshell = 10\;\si{m}$:
\begin{equation}
\Delta a_{\rm tidal} = \frac{2 \times 10^3}{10}\times 5
= 1000\;\si{m/s^2} \approx 100g.
\end{equation}
This is significant but not necessarily destructive for a solid
metallic object. For a human body ($L = 2\;\si{m}$):
\begin{equation}
\Delta a_{\rm tidal} = \frac{2\times 10^3}{10}\times 2
= 400\;\si{m/s^2} \approx 40g.
\end{equation}
This would be injurious during transit through the boundary, but
the transit time is:
\begin{equation}
\Delta t_{\rm transit} \sim \frac{\Lshell}{v}
\sim \frac{10}{200} = 0.05\;\si{s}.
\end{equation}
A 40$g$ impulse for 50~ms is survivable (comparable to automotive
crash loading).

\paragraph{Key design principle.}
Tidal forces scale as $L/\Lshell^2$. A thicker shell (larger
$\Lshell$) \emph{dramatically} reduces tidal stresses---this is
the ``gentle capture'' regime. A shell thickness of
$\Lshell \sim 100\;\si{m}$ would reduce tidal forces by a factor
of 100 relative to $\Lshell = 10\;\si{m}$.

\paragraph{Interior tidal forces.}
Inside the bubble ($r < \Rb - \Lshell$), the $\psi$-field is
approximately constant. The tidal forces are negligible:
$\Delta a_{\rm tidal} \to 0$. Objects inside the bubble float in
perfect freefall.


%=============================================================================
\section{EM Signal Propagation Across the Bubble}\label{sec:em}
%=============================================================================

\subsection{Refractive Index Distribution}

The effective refractive index inside and around the bubble is
(from DFD, $n = e^{\psi}$):
\begin{equation}
n(\mathbf{x}) = \exp\!\left[\psi_{\rm bg}(\mathbf{x})
  + \psi_{\rm eng}(\mathbf{x})\right].
\end{equation}
In the shell region, the refractive index deviates from unity by:
\begin{equation}
\delta n = n - 1 \approx \psi_{\rm eng}
\quad (\text{for } |\psi_{\rm eng}| \ll 1).
\end{equation}
For our parameters, $\delta n \sim 10^{-13}$---an extraordinarily
small perturbation.

\subsection{Ray Deflection (Gravitational Lensing)}

A light ray passing through the bubble shell experiences a
deflection angle (from Fermat's principle, DFD Eq.~2.3):
\begin{equation}
\theta_{\rm defl} = -\int \frac{\nabla_\perp n}{n}\,ds
\approx -\int \nabla_\perp\psi_{\rm eng}\,ds.
\end{equation}
For a ray passing through the shell at impact parameter $b$
from the bubble centre:
\begin{equation}
\theta_{\rm defl} \sim \frac{2\psib\Lshell}{b}
\sim \frac{4a_{\rm max}\Lshell^2}{c^2\,b}.
\end{equation}

\paragraph{Numerical estimate.}
For $a_{\rm max} = 100g$, $\Lshell = 10\;\si{m}$, $b = 50\;\si{m}$:
\begin{equation}
\theta_{\rm defl} \sim \frac{4\times 10^3\times 100}
  {9\times 10^{16}\times 50} \sim 10^{-13}\;\text{rad}
\approx 0.02\;\text{nanoarcseconds}.
\end{equation}
This is \textbf{undetectable} by any current or foreseeable
optical instrument.

\subsection{Spectral Shift of EM Signals}

A photon climbing out of the $\psi$-well experiences a
gravitational redshift (DFD Eq.~2.5):
\begin{equation}
\frac{\Delta\lambda}{\lambda} = \frac{\Delta\nu}{\nu}
= \Delta\psi = \psi_{\rm interior} - \psi_{\rm exterior}
\approx \psib.
\end{equation}
For $\psib \sim 10^{-13}$, this is a fractional frequency shift
of $10^{-13}$---detectable only by atomic clocks or precision
spectroscopy, not by radar or IR sensors.

\subsection{Phase Delay (Shapiro-like Effect)}

A radar signal traversing the bubble experiences a round-trip
time delay:
\begin{equation}
\Delta t = \frac{2}{c}\int \psi_{\rm eng}\,ds
\sim \frac{2\psib\Lshell}{c}.
\end{equation}
Numerically: $\Delta t \sim 2\times 10^{-13}\times 10\,/\,
(3\times 10^8) \sim 10^{-21}\;\si{s}$. This is
$\sim 10^{-6}$ of a radar wavelength at $10\;\si{GHz}$---
utterly negligible.

\subsection{Radar Return from an Object Inside the Bubble}

A radar pulse enters the bubble, traverses the shell
(negligible deflection and delay), illuminates the object, and
returns. The radar cross-section of the captured object is
essentially \textbf{unchanged}. The bubble is transparent to EM
radiation at these $\psi$-values.

\paragraph{Conclusion.}
At the $\psi$-values needed for $100g$-class acceleration
($\psib \sim 10^{-13}$), the bubble is \textbf{electromagnetically
invisible}. It produces no detectable lensing, spectral shift,
or radar anomaly.


%=============================================================================
\section{External Observer Phenomenology}\label{sec:observables}
%=============================================================================

\subsection{What External Observers See}

\paragraph{Visual appearance.}
The bubble and its contents are optically transparent (to the
precision of $\delta n \sim 10^{-13}$). An observer looking at the
bubble sees the ship and any captured objects with no detectable
distortion, colour shift, or magnification.

\paragraph{Kinematic anomaly.}
The \textbf{only} observable signature is kinematic: objects that
enter the bubble's gradient zone suddenly co-accelerate with the
bubble. A missile that was tracking on a ballistic trajectory
suddenly matches the bubble's velocity and trajectory.

From the external frame, this appears as:
\begin{enumerate}
\item The missile enters a zone near the ship.
\item The missile's velocity rapidly changes to match the ship's.
\item The missile appears to ``stick'' to the ship and co-move.
\end{enumerate}
This is indistinguishable from the missile being physically grabbed,
except there is no impact, no debris, and no deceleration of the ship.

\paragraph{Apparent luminosity.}
The ship's luminosity (thermal emission, reflected light) is
unaffected to measurable precision. No dimming or brightening
accompanies the $\psi$-bubble.

\subsection{What the Ship's Occupants Experience}

\paragraph{Zero-g environment.}
Inside the bubble, the $\psi$-field is uniform. All objects---crew,
instruments, the captured missile---are in freefall. The occupants
experience \textbf{weightlessness}, exactly like astronauts in
orbit.

\paragraph{The captured missile floats.}
From the interior perspective, the missile simply drifted in from
outside (with some residual relative velocity from tidal effects
during shell transit) and now floats nearby. There is no sensation
of ``capture''---just an object arriving and then being weightless.


%=============================================================================
\section{Atmospheric Transit}\label{sec:atmosphere}
%=============================================================================

\subsection{Air Molecules Entering the Gradient}

When the bubble moves through atmosphere, air molecules at the
leading edge encounter the $\psi$-gradient. By the same freefall
mechanism (Eq.~\ref{eq:eom}), each air molecule accelerates at
$(c^2/2)\nabla\psi$---the same acceleration as the ship.

\subsection{Molecular Dynamics at the Shell Boundary}

Consider the leading shell of the bubble moving at velocity $v$
through air of density $\rho_{\rm air}$:

\begin{enumerate}
\item \textbf{Gradient encounter}: An air molecule at rest (in the
ground frame) enters the gradient zone. It begins to accelerate
in the direction of the bubble's motion.

\item \textbf{Acceleration through the shell}: Traversing the shell
thickness $\Lshell$, the molecule gains velocity:
\begin{equation}
\Delta v_{\rm air} \sim a_{\rm max}\,\frac{\Lshell}{v}
= \frac{c^2\psib}{2v}.
\end{equation}

\item \textbf{If $\Delta v_{\rm air} \geq v$}: The molecule is fully
captured and entrained. It co-moves with the bubble. The air
\emph{inside} the bubble is carried along.

\item \textbf{If $\Delta v_{\rm air} < v$}: The molecule is partially
deflected but passes through. This generates a partial
entrainment and possible heating.
\end{enumerate}

\subsection{Full Entrainment Condition}

Full air entrainment requires:
\begin{equation}
\frac{c^2\psib}{2v} \geq v
\quad\Longrightarrow\quad
v \leq c\sqrt{\frac{\psib}{2}} = v_{\rm capture}.
\end{equation}
This is the same capture velocity derived in
Section~\ref{sec:missile}. For $\psib \sim 10^{-13}$,
$v_{\rm capture} \sim 200\;\si{m/s}$.

\subsection{Supersonic Regime}

For $v > v_{\rm capture}$ (e.g., supersonic or hypersonic flight),
air molecules are \emph{not} fully captured. They pass through the
gradient zone with partial deflection. However, the energy transfer
from the bubble to the air is:
\begin{equation}
\Delta E_{\rm per\,molecule} = \frac{1}{2}m_{\rm air}(\Delta v)^2
= \frac{m_{\rm air}\,c^4\psib^2}{8v^2}.
\end{equation}
The power dissipated per unit area of bubble cross-section:
\begin{equation}
P/A = n_{\rm air}\,v\,\Delta E_{\rm per\,molecule}
= \frac{n_{\rm air}\,m_{\rm air}\,c^4\psib^2}{8v},
\end{equation}
where $n_{\rm air}$ is the number density. This is
\emph{extraordinarily} small for $\psib \sim 10^{-13}$.

\paragraph{Implication.}
At low $\psi$-values ($\psib \sim 10^{-13}$), the bubble is
essentially invisible to the atmosphere at supersonic speeds.
Air passes through the gradient with negligible deflection.
There is no shock wave, no sonic boom, and no thermal signature
from aerodynamic heating.

For high-$\psi$ bubbles ($\psib \sim 10^{-9}$ or above), full
atmospheric entrainment becomes possible at high speeds, and
the bubble would carry a co-moving air column, effectively
eliminating aerodynamic drag.

\subsection{The Two-Regime Summary}

\begin{table}[h]
\caption{Atmospheric interaction regimes.}
\begin{ruledtabular}
\begin{tabular}{lcc}
Regime & Condition & Behavior \\
\hline
Transparent & $v \gg v_{\rm capture}$ &
  Air passes through; no drag, no shock \\
Entraining & $v \lesssim v_{\rm capture}$ &
  Air captured; co-moving atmosphere \\
\end{tabular}
\end{ruledtabular}
\end{table}

In the entraining regime, the bubble carries its own atmosphere.
Objects within (including the ship and occupants) experience a
pressurized environment. In the transparent regime, the interior
must be independently pressurized.


%=============================================================================
\section{High-$\psi$ Regime: Strongly Interacting Bubbles}\label{sec:high-psi}
%=============================================================================

\subsection{Scaling to Higher Accelerations}

For applications requiring capture of fast-moving objects
(e.g., hypersonic missiles at Mach~10, $\sim 3.4\;\si{km/s}$):
\begin{equation}
\psib = \frac{2v_{\rm capture}^2}{c^2}
= \frac{2\times (3400)^2}{9\times 10^{16}}
\sim 2.6\times 10^{-10}.
\end{equation}
This requires $a_{\rm max} \sim 10^6\;\si{m/s^2} \sim 10^5\,g$
for $\Lshell = 10\;\si{m}$, or $\sim 10^3\,g$ for
$\Lshell = 1\;\si{km}$.

At this $\psib$-value:
\begin{itemize}
\item Spectral shift: $\Delta\lambda/\lambda \sim 10^{-10}$
  (detectable by atomic clocks but not optical instruments).
\item Lensing: $\theta \sim 10^{-10}\;\text{rad}
  \sim 20\;\mu\text{arcsec}$ (at the edge of microarcsecond
  astrometry).
\item Phase delay: $\Delta t \sim 10^{-18}\;\si{s}$
  (detectable by femtosecond lasers).
\end{itemize}

\subsection{Observable Signatures at High $\psi$}

Only at $\psib \gtrsim 10^{-6}$ (requiring acceleration
$\sim 10^{10}\,g$ at $\Lshell = 10\;\si{m}$) would the bubble
produce readily detectable optical signatures:
\begin{itemize}
\item Gravitational lensing distortion visible to telescopes
\item Measurable spectral shifts of transmitted light
\item Radar timing anomalies detectable by standard equipment
\end{itemize}

This regime corresponds to accelerations comparable to the surface
gravity of a neutron star. The energy requirements would be
correspondingly extreme.


%=============================================================================
\section{Formal Summary}\label{sec:summary}
%=============================================================================

\subsection{Principal Results}

All results derived from DFD Eq.~2.14:
$\mathbf{a} = (c^2/2)\nabla\psi$ and the optical metric
$n = e^{\psi}$.

\begin{enumerate}
\item \textbf{Capture radius}
(Eq.~\ref{eq:capture-radius}):
\begin{equation*}
\Rc = \Rb + \frac{\Lshell}{2}\ln\!\left(
  \frac{4\,a_{\rm max}}{a_{\rm thresh}}\right).
\end{equation*}
Logarithmic extension beyond the shell; controlled by
$\Lshell$.

\item \textbf{Tidal force at boundary}
(Eq.~\ref{eq:tidal-max}):
\begin{equation*}
\Delta a_{\rm tidal} = \frac{2\,a_{\rm max}}{\Lshell}\,L.
\end{equation*}
Survivable for thick shells ($\Lshell \gg L$).

\item \textbf{EM transparency}: At propulsion-scale
$\psi$-values ($\psib \sim 10^{-13}$), the bubble produces
\emph{no detectable} lensing ($\theta \sim 10^{-13}\;\text{rad}$),
spectral shift ($\Delta\lambda/\lambda \sim 10^{-13}$),
or radar anomaly ($\Delta t \sim 10^{-21}\;\si{s}$).

\item \textbf{Atmospheric transparency}: For $v > v_{\rm capture}
= c\sqrt{\psib/2}$, air passes through the bubble undeflected.
No shock wave, no sonic boom.

\item \textbf{Missile capture by freefall}: Any object entering
the gradient zone accelerates identically to the ship (WEP).
Inside the bubble, the relative acceleration reduces to tidal
terms only. The object co-moves in freefall.
\end{enumerate}

\subsection{The Core Physical Insight}

The $\psi$-bubble does not ``push'' or ``pull'' objects with a
force. It creates a region of modified spacetime refractive index
in which all objects---ship, crew, missiles, air molecules---fall
freely. Capture is not mechanical; it is gravitational.

From the DFD perspective, the bubble generates a local gravity well
that moves through space. Everything within the well's gradient
falls with it. This is no different in principle from the Moon
being ``captured'' by Earth's gravity---except that the well is
artificial, mobile, and switchable.

\end{document}
